\section{The homological eigenvalue}\label{sec:homology}

The main coordinate of \(f\) is a bounded additive perturbation of
squaring. Fading out that perturbation gives the homological information
needed for Theorem~\ref{thm:main}.

\begin{proposition}\label{prop:homology}
For the map \(f\) in \eqref{eq:map}, the action
\(f_*:H_2(M;\R)\to H_2(M;\R)\) has eigenvalue \(2\).
\end{proposition}

\begin{proof}
Let \(\pi:M\to\Sigma\) be projection to \(z\), and put \(g(z)=z^2\),
with \(g(\infty)=\infty\). Define
\[
 C(t,v)=20\sum_{\ell=1}^q\bar\beta_\ell(t)S(v_\ell).
\]
Since \(|S|\le3\), we have \(|C|\le60q\). For \(0\le s\le1\), set
\[
 H_s(t,z,v)=z^2+sC(t,v)
\]
when \(z\) is finite, and \(H_s(t,\infty,v)=\infty\).
This is a continuous homotopy from \(g\pi\) to \(\pi f\).
Indeed, in the inverse coordinate \(w=1/z\) at infinity, the output
has inverse coordinate
\[
 \frac{w^2}{1+sC(t,v)w^2}.
\]
For all \(s,t,v\), the denominator is nonzero on one fixed
neighborhood of \(w=0\), since \(|sC|\le60q\).
This verifies continuity jointly in every variable, including \(s\).

By homotopy invariance and functoriality of singular homology
\cite[Theorem~2.10]{Hatcher2002},
\begin{equation}\label{eq:intertwining}
 \pi_*f_*=g_*\pi_*
 \quad\text{on }H_2(M;\R).
\end{equation}
The oriented sphere has \(H_2(\Sigma;\R)\cong\R\), and the action
of \(g\) on this group is multiplication by its degree.
The regular value \(1\) has the two preimages \(1\) and \(-1\).
At either preimage the derivative is complex multiplication by
\(2z\), with positive real determinant \(|2z|^2=4\).
Both local degrees are \(+1\), so \(\deg g=2\)
\cite[Proposition~2.30 and Lemma~2.49]{Hatcher2002}.

The projection \(\pi\) has a section obtained by fixing the clock
and all register coordinates. Therefore
\(\pi_*:H_2(M;\R)\to H_2(\Sigma;\R)\) is surjective.
Write \(A=f_*|_{H_2(M;\R)}\) and \(P=\pi_*\). Equation
\eqref{eq:intertwining} gives \(P(A-2I)=0\), with \(P\ne0\).
If \(A-2I\) were invertible, this equality would imply \(P=0\).
Thus \(2\) is an eigenvalue of \(A\).
\end{proof}

\begin{proof}[Proof of Theorem~\ref{thm:main}]
Proposition~\ref{prop:regularity} constructs a \(C^1\) self-map of
the stated compact smooth manifold.
Proposition~\ref{prop:zero-entropy} gives \(\htop(f)=0\), and
Proposition~\ref{prop:homology} supplies the eigenvalue \(2\) on
\(H_2(M;\R)\). This group is an invariant summand of \(H_*(M;\R)\),
so \(\rho(f_*)\ge2\). Finally, \(S\) vanishes on a neighborhood
of infinity. Holding all other coordinates fixed and varying
one register within that neighborhood leaves every output in
\eqref{eq:map} unchanged. Hence \(f\) is noninvertible.
\end{proof}
